\documentclass[12pt]{article}
\usepackage[T1]{fontenc}
\usepackage[utf8]{inputenc}
\usepackage{lmodern}
\usepackage{textcomp}
\usepackage{enumitem}
\usepackage{geometry}
\usepackage{graphicx}
\usepackage{amsmath, amssymb}
\geometry{margin=1in}
\begin{document}
\begin{center}
Astronomy - Graduate Level Assessment 20 \
Total Marks: 40 \
Answer all questions.
\end{center}

\section*{Multiple Choice (10 marks)}
\begin{enumerate}[label=\arabic*.]
\item The lunar maria are:
\begin{enumerate}[label=(\alph*)]
    \item ancient heavily cratered highlands
    \item dark lavas inside volcanic calderas
    \item dark lavas filling older impact basins
    \item the bright regions on the Moon
\end{enumerate}
\item Why are Cepheid stars relevant for astronomers?
\begin{enumerate}[label=(\alph*)]
    \item To measure interstellar mass.
    \item To measure galactic distances.
    \item To measure galactic energy-density.
    \item To measure interstellar density.
\end{enumerate}
\item The sky is blue because
\begin{enumerate}[label=(\alph*)]
    \item the Sun mainly emits blue light.
    \item the atmosphere absorbs mostly blue light.
    \item molecules scatter red light more effectively than blue light.
    \item molecules scatter blue light more effectively than red light.
\end{enumerate}
\item The terrestrial planet cores contain mostly metal because
\begin{enumerate}[label=(\alph*)]
    \item the entire planets are made mostly of metal.
    \item metals condensed first in the solar nebula and the rocks then accreted around them.
    \item metals sank to the center during a time when the interiors were molten throughout.
    \item radioactivity created metals in the core from the decay of uranium.
\end{enumerate}
\item Which one of these constellations is not located along the Milky Way in the sky?
\begin{enumerate}[label=(\alph*)]
    \item Perseus
    \item Cygnus
    \item Scorpius
    \item Leo
\end{enumerate}
\end{enumerate}

\section*{True / False (10 marks)}
\textit{Answer True or False}
\begin{enumerate}[label=\arabic*.]
\setcounter{enumi}{5}
\item True or False: Knowing both the actual brightness and apparent brightness of an object allows you to determine its composition.
\item True or False: The ratio of solar radiation flux on Mercury's surface at perihelion (0.304 AU) compared to aphelion (0.456 AU) is 9:4.
\item True or False: Differentiation occurs solely due to spontaneous emission from radioactive atoms in a planetary interior.
\item True or False: The visible part of the electromagnetic spectrum is typically defined as spanning from 240 to 680 nanometers.
\item True or False: The part closer to the axis of a solid disk rotates with a longer period than the outer regions of the disk.
\end{enumerate}

\section*{Long Form Response (20 marks)}
\textit{Answer each question with a clear and complete explanation.}
\begin{enumerate}[label=\arabic*.]
\setcounter{enumi}{10}
\item Discuss the significance of the Laniakea supercluster in the context of the Milky Way's position within the universe. Analyze its impact on our understanding of cosmic structure.
\item Discuss the definition of a parsec in astronomy, including its significance and the method by which it is determined in relation to astronomical units and angles.
\end{enumerate}

\vfill
\noindent\textit{End of Paper}
\end{document}